% =====================================================================
% Section 4 body: Benchmark and Mechanism-Ablation Methodology.
% Master supplies the \section header. This fragment is the BODY only.
% Organized to LEAD with external (non-in-sample) component benchmarks, then the
% mechanism-attribution (ablation) experiment, then the in-sample defense.
% NOTE: the "component intercomparison" here benchmarks closure BRANCHES against
% base-pressure datasets; it does NOT swap closures at apogee (out of scope).
% All numbers per paper/_shared/CANONICAL_FACTS.md sections B, C, D.
% Citation rule (Sec. H): the empirical Cd_base = 0.064 + 0.186/M^2 is an
% EMPIRICAL fit validated vs NACA TN 3393 and consistent with ESDU 77021;
% NO Devan-Ashwood. Companion JSR paper = \cite{yu2026jsr} (placeholder).
% =====================================================================

The methodology has three parts, ordered from the most external evidence to
the most in-sample. The \emph{component intercomparison}
(Section~\ref{sec:method-component}) is the substantive physics of the paper: it
asks how each candidate base-drag closure reproduces \emph{measured} base
pressure in isolation, against external wind-tunnel and free-flight data that
played no part in any corpus calibration. The \emph{mechanism ablation}
(Section~\ref{sec:method-ablation}) is the motivating experiment: it asks how
much each implemented aerodynamic mechanism actually moves integrated apogee, by
disabling it one at a time across the flight corpus, and it is what singles out
base drag as the dominant driver. The \emph{decontaminated holdout}
(Section~\ref{sec:method-holdout}) is the in-sample defense: it asks whether the
two corpus-frozen base-drag scale constants generalize to flights that never
informed their calibration, pre-empting the circularity critique that the
ablation substrate would otherwise invite. We lead with the externally validated
closure fidelity, use the ablation as the finding that motivates a base-drag
study at all, and confront the in-sample constants head-on rather than letting
them read as the result.

\subsection{Component intercomparison protocol}
\label{sec:method-component}

Each base-drag closure introduced in Section~\ref{sec:basedrag-models} is
exercised against published base-pressure data with the boundary-layer state and
length Reynolds number $\mathrm{Re}_{L}$ matched to the reference condition, so
that closures are compared on the flow regime each was derived for rather than
averaged across regimes. The discriminator is deliberate: as
Section~\ref{subsec:laminar} shows, applying a turbulent correlation to laminar
separation data inflates the error by an order of magnitude, so the protocol
conditions closure selection on the boundary-layer state and not on Mach number
alone. The error metric throughout is the mean absolute percentage error (MAPE)
of the predicted base-pressure or base-drag coefficient against the digitized
reference points. Three external datasets anchor the comparison, none of which
was used to tune any integrated-trajectory constant; a fourth, computational,
comparator is identified but deferred.

\paragraph{NACA TN 3393 (turbulent and laminar).} The primary anchor is the
Reller and Hamaker base-pressure investigation of nonlifting bodies of revolution
at Mach 2.73--4.98~\cite{chapman1955}, which reports base pressure for both
turbulent and laminar boundary-layer states on the same geometries and therefore
exercises both base-drag branches against a single, internally consistent
experiment. Against the turbulent points the empirical correlation
$C_{d,\mathrm{base}} = 0.064 + 0.186/M^{2}$ of Eq.~\eqref{eq:empirical} attains a
MAPE of $15.9\%$; against the laminar points the Chapman laminar closure of
Eq.~\eqref{eq:laminar} attains $4.4\%$~\cite{chapman1955}. The laminar closure is
thus markedly more faithful on its target regime, and the two branches are not
interchangeable: applying the turbulent correlation to the same laminar data
yields a MAPE of $44\%$. This is the quantitative basis for treating ``the''
base-drag closure as a state-dependent family rather than a single object.

\paragraph{Hart NACA RM L52E06 (transonic turbulent).} To extend the turbulent
comparison below the TN~3393 floor and through the transonic base-drag peak near
$M \approx 1.05$, the turbulent closure is additionally compared against the Hart
free-flight base-pressure measurements (NACA~RM~L52E06~\cite{hart1952basedrag}),
where it attains a MAPE of $4.0\%$. The two turbulent anchors together span the transonic
peak (Hart) and the supersonic decay (TN~3393), and bracket the band in which the
empirical correlation $C_{d,\mathrm{base}} = 0.064 + 0.186/M^{2}$ and the degree-4
transonic blend of Section~\ref{subsec:transonic} operate. The Hart finless
free-flight reading of $0.250 \pm 0.013$ near the peak also fixes the transonic
anchor point of that blend.

\paragraph{Chapman laminar data.} The laminar branch is anchored on the laminar
subset of the same TN~3393 experiment~\cite{chapman1955}, against which the
Chapman laminar closure attains the $4.4\%$ MAPE quoted above. This isolates the
``perfect-finish'' / delayed-transition case relevant to highly polished slender
vehicles, for which the turbulent correlation over-predicts base drag, and
provides the externally validated counterpart to the laminar branch invoked for
the smoothest airframes in the corpus.

\paragraph{Transonic CFD comparator (deferred, disclosed).} A fourth comparison is
identified but not executed in this work: the Sahu, Nietubicz and Steger
thin-layer Navier--Stokes computation of base flow on a secant-ogive-cylinder-%
boattail projectile at transonic speed
($M \approx 0.9$--$1.2$)~\cite{sahu1983} would provide an independent, CFD-grade
check of the closures through the transonic peak, where the empirical and
Chapman--Korst forms hand off to the degree-4 blend. Its geometry is structurally
different from the projectile families used for the other anchors, so reproducing
it would require building a separate vehicle model and digitizing the published
base-pressure distributions. We disclose this comparison as \emph{deferred}---it
is documented as a memo only and is not exercised as a quantitative comparator
here---rather than report a partial result. We claim \emph{no computational fluid
dynamics of our own}; \cite{sahu1983} is a published third-party computation cited
as a prospective comparator.

\subsection{Controlled mechanism-ablation protocol}
\label{sec:method-ablation}

The component intercomparison establishes closure fidelity in isolation but does
not establish which mechanism governs the \emph{integrated} apogee error. A
closure that matches a wind-tunnel point to a few percent may still move apogee by
much more, or much less, than its component error suggests, because apogee
integrates the drag history over a trajectory weighted toward the lower-Mach
portion of the flight. To resolve this, we run a controlled ablation on the
archived 24-flight mechanism-ablation subset of the externally selected
supersonic altitude-comparison corpus, spanning Mach~$0.54$--$3.46$, used for
the integrated validation reported in the companion flagship
study~\cite{yu2026jsr}. The ablation subset contains $23$ single-stage flights
plus the AeroPac~104K two-stage closure. MESOS~293K remains part of the
companion 25-flight validation corpus and is reported there as the highest-Mach
two-stage case; it is not part of the archived mechanism-ablation artifact from
which the 8.10/0.87/0.39/0.15/0.00~pp ranking is computed. Here the corpus is
the \emph{substrate for a mechanism-attribution experiment}, not the basis for
an accuracy headline.

The protocol disables one aerodynamic mechanism at a time, re-flies all $24$
of these trajectories, and records the change in the corpus mean absolute apogee error,
$\overline{|\Delta|}$, relative to the fully enabled baseline. Five mechanisms are
ablated: the finned-base drag augmentation $\Phi_{\text{fin}}$ of
Eq.~\eqref{eq:finnedbase}; the Van~Driest~II compressible skin-friction
transformation~\cite{hopkins1971}; the DATCOM Section~4.1.5.1 fin wave-drag
closure~\cite{datcom1978,ulmann1956}; the shock-geometry pre-pass; and the
Pitts--Nielsen--Kaattari fin--body interference correction, together with its
$K_{1}$ floor.

Each mechanism is ablated \emph{in isolation}---disabled singly, with every other
mechanism held at its baseline state, and the corpus re-run from a clean
configuration for each ablation. This isolation discipline is essential and not
incidental: the ablation flags are evaluated as shared static state, so toggling
several mechanisms within a single process risks cross-contaminating the recorded
baseline and corrupting the attribution. Running each ablation as an independent,
clean corpus pass removes that contamination, and the resulting single-mechanism
apogee sensitivities are the per-mechanism $\overline{|\Delta|}$ values reported
in Section~\ref{sec:results}.

The ranking that emerges---and the finding that motivates this paper---is steeply
ordered. The finned-base augmentation dominates the apogee error budget at a mean
$\overline{|\Delta|}$ of $8.10$~pp (reaching $39.5$~pp on the Mach-2.19 A-601
Kinsel flight), more than an order of magnitude larger than the next mechanism.
Compressible skin friction contributes $0.87$~pp, the DATCOM fin wave-drag term
$0.39$~pp, and the shock-geometry pre-pass---the celebrated architectural
innovation of the companion paper---only $0.15$~pp, because that pre-pass is inert
below Mach~1 and apogee integrates a trajectory dominated by lower-Mach drag; the
PNK interference term, with its $K_{1}$ floor, registers no measurable apogee
effect ($0.00$~pp). Base drag, not shock geometry, governs the apogee error
budget, which is what licenses a focused base-drag closure study distinct from the
integrated parity result of~\cite{yu2026jsr}.

\subsection{Decontaminated prospective holdout for the corpus-frozen constants}
\label{sec:method-holdout}

The turbulent base-drag branch carries two scale constants that were frozen
against the corpus rather than against an independent component dataset: a
thick-boundary-layer scale, \texttt{THICK\_BL\_K}~$=2.2$, and a slender-body
scale, \texttt{SLENDER\_BODY\_K}~$=0.0025$. Because these constants were tuned on
corpus flights, the integrated result is \emph{partly in-sample}, and a naive
corpus-wide accuracy figure would be optimistically biased. We disclose this
plainly and test generalization with a decontaminated prospective holdout, which
is the defense against the circularity critique that follows from using a
corpus-calibrated quantity as the ablation substrate of
Section~\ref{sec:method-ablation}.

The decontamination rule is strict: \emph{every} flight that either constant
touched during calibration---the anchor flights for both constants and the
source-diagnostic flight---is quarantined in the development partition. This
decontamination is applied as an overlay on the corpus's existing
development/holdout bookkeeping: the five flights involved in fixing the two
constants (Raven, Rabia, Rabia Short Fin Can, Kinsel, and Torrent) are placed in
the development partition (DEV) together with the cases already reserved for
closure development and ablation, while no calibration-touched flight is allowed
to remain in the blind set (HOLDOUT). This yields a development partition of
thirteen flights and a holdout of twelve. We stress that this calibration
decontamination is a stricter criterion than, and is not to be conflated with,
the corpus's prospective forward-freeze protocol, which partitions flights by
whether they may be re-tuned after a freeze date rather than by whether they
informed the two scale constants, and therefore assigns three of the
calibration-touched flights differently. The constants are frozen on DEV and
never re-touched; HOLDOUT is scored exactly once. If the constants merely overfit
the corpus, HOLDOUT accuracy should degrade relative to DEV; if they capture a
transferable physical scaling, HOLDOUT accuracy should hold or improve.

The split yields a development partition of $n=13$ (mean signed error $+0.22\%$,
MAE $5.47\%$) and a blind holdout of $n=12$ (mean signed error $-1.03\%$, MAE
$3.95\%$). These partition statistics are computed from the per-flight apogee
errors of the archived corpus-baseline artifact, evaluated against the
current-code MESOS~293K value, so the split is regenerable from the released data
under the membership stated above. The holdout is \emph{more} accurate than the
development partition ($3.95\% < 5.47\%$), which is the primary in-sample defense:
the two corpus-frozen base-drag constants generalize to unseen flights rather than
overfitting the calibration set. This result is what permits the corpus to serve
as the ablation substrate in Section~\ref{sec:method-ablation} without
circularity, and it is consistent with the holdout reported in the companion
study~\cite{yu2026jsr}.
